\begin{pastexam}[label={px:ntu-107-dsa-06}]{臺大資工所 107 學年度 第 VI 題}
（問答）

(25 points) The network of roads in NTU can be modeled as a graph $G=(V,E)$, where each node
$v_i\in V$ represents a building and each edge $e_{ij}\in E$ represents a road connecting the
buildings $v_i$ and $v_j$. Each edge $e_{ij}$ is associated with a positive weight $w_{ij}$, i.e.,
$w_{ij}>0$.
\begin{enumerate}[label=(\alph*)]
\item (10 points) Consider all roads are bidirectional. Before an earthquake, it is possible to
travel between any pair of buildings. However, after the earthquake, all of the roads are damaged
and closed until repairs are completed. Suppose the weight $w_{ij}$ indicates the time required to
repair the corresponding road $e_{ij}$. Assume that all roads are fixed at the same time.

Given two buildings $v_a$ and $v_b$, please propose an efficient algorithm to determine a path
starting from $v_a$ to $v_b$ that can be fixed in the smallest amount of time. Analyze the time
complexity of your algorithm.
\item (15 points) Consider all roads are directional but we don't know their directions in $G$.
Suppose the weight $w_{ij}$ indicates the time required to travel from the building $v_i$ to the
building $v_j$. Given two buildings $v_a$ and $v_b$ and the following two different conditions,
for each condition, please propose an efficient algorithm to determine the directions of the roads
in $G$ so that the time to travel from $v_a$ to $v_b$ is minimum. Assume that $v_b$ is reachable
from $v_a$ in $G$.
\begin{enumerate}[label=(\arabic*)]
\item (5 points) Consider all edges have equal weights. Your algorithm should run in linear time.
\item (10 points) Consider the edges might have different weights. Analyze the time complexity of
your algorithm in this question.
\end{enumerate}
\end{enumerate}
\end{pastexam}

\begin{pastsoln}
本題 (a) 是最小生成樹的瓶頸性質（資結書第~24~章），(b) 是最短路徑；(b) 佔的分數較多，歸在最短路徑
這一章。 以下 $n = |V|$、$m = |E|$。

\textbf{(a) 題意。} 所有道路同時開始修，道路 $e_{ij}$ 在時間 $w_{ij}$ 修好。 一條從 $v_a$ 到 $v_b$ 的路徑
要等它上面\textbf{每一段}都修好才能走，所以它能通行的時間是路徑上最大的 $w_{ij}$，也就是路徑的瓶頸
（資結書第~24.5.3~節）。 題目要的是瓶頸最小的路徑，不是權重總和最小的最短路徑。

\textbf{演算法。} 照 Kruskal 演算法由輕到重加邊，$v_a$、$v_b$ 第一次連通時停下來，在選到的邊組成的
森林上找 $v_a$ 到 $v_b$ 的路徑。
\begin{pseudo}{最早修好的路徑}
\Function{Earliest-Path}{$G, w, a, b$}
    \State $F \gets \emptyset$
    \For{每個頂點 $v \in V$}
        \State \Call{Make-Set}{$v$}
    \EndFor
    \State 把 $E$ 的邊依權重由小到大排序
    \For{依排好的順序，每條邊 $(u, v)$}
        \If{$\Call{Find-Set}{a} = \Call{Find-Set}{b}$}
            \State \Break
        \EndIf
        \If{$\Call{Find-Set}{u} \ne \Call{Find-Set}{v}$}
            \State $F \gets F \cup \{(u, v)\}$；\Call{Union}{$u, v$}
        \EndIf
    \EndFor
    \State 在森林 $F$ 上從 $a$ 做 BFS，沿 $\pi$ 從 $b$ 走回 $a$，得到路徑 $P$
    \State \Return $P$
\EndFunction
\end{pseudo}
地震前任兩棟建築都連通，所以 $G$ 連通，迴圈一定會在某條邊之後讓 $a$、$b$ 連通（$a = b$ 時一開始就
停，路徑只有一個頂點）。 $F$ 是森林，$a$、$b$ 在同一棵樹裡，樹上兩點之間恰好有一條簡單路徑，第 11 行
找到的就是它。

\textbf{正確性。} 設 $e^{*}$ 是最後一條加進 $F$ 的邊，也就是讓 $a$、$b$ 第一次連通的那一條，
$w^{*} = w(e^{*})$。
\begin{itemize}
\item \textbf{$P$ 的瓶頸 $\le w^{*}$。} $P$ 的邊都在 $F$ 裡，$F$ 的邊都不比 $e^{*}$ 晚處理，權重都
$\le w^{*}$。
\item \textbf{任何路徑的瓶頸 $\ge w^{*}$。} 令 $E^{-}$ 是排序後排在 $e^{*}$ 前面的邊。 處理完 $E^{-}$ 時，
互斥集合正好是圖 $(V, E^{-})$ 的連通元件：被略過的邊兩端本來就在同一個集合，加了也不改變連通性。
那時 $a$、$b$ 還在不同的集合，所以只用 $E^{-}$ 的邊從 $a$ 走不到 $b$。 任何一條 $a$ 到 $b$ 的路徑 $Q$
都至少用到一條不在 $E^{-}$ 的邊，它排在 $e^{*}$ 或它後面，權重 $\ge w^{*}$，所以 $Q$ 的瓶頸
$\ge w^{*}$。
\end{itemize}
兩點合起來，$P$ 的瓶頸是所有路徑裡最小的，它在時間 $w^{*}$ 就能通行，沒有路徑能更早。 換個說法：
$F$ 是 Kruskal 求最小生成樹的前半段，$P$ 是最小生成樹上 $a$ 到 $b$ 的路徑，由資結書定理~24.13，
它的瓶頸最小。

\textbf{時間。} 排序 $O(m \lg m) = O(m \lg n)$；最多 $m$ 輪迴圈，每輪幾次 \proc{Find-Set}，用依大小合併
加路徑壓縮共 $O(m\,\alpha(n))$（資結書第~13~章）；$F$ 最多 $n - 1$ 條邊，BFS 是 $O(n)$。 合計
$O(m \lg n)$，瓶頸在排序。

\textbf{另一種寫法。} 把 Dijkstra 演算法的「$d[u] + w(u, v)$」改成「$\max(d[u], w(u, v))$」，$d[v]$ 代表
目前找到的最小瓶頸，取出 $d$ 最小的頂點時它就確定了，證明和 Dijkstra 相同（$\max$ 和加非負數一樣，
沿路徑不會變小）。 用二元堆積也是 $O(m \lg n)$，而且找到 $v_b$ 就能停。 兩種寫法都要說明的關鍵是：
題目要的是\textbf{瓶頸}，直接跑一般的最短路徑（權重相加）會答錯。 例如 $v_a$ 到 $v_b$ 一條路是一段
$10$，另一條路是三段 $4$：總和最小的是 $10$ 那條，但三段 $4$ 的那條在時間 $4$ 就修好了。

\textbf{(b) 題意。} 每條路只能單向通行，方向由我們決定；$w_{ij}$ 是走這條路的時間（不論定成哪個方向）。
要把每條路定一個方向，讓定完之後 $v_a$ 到 $v_b$ 的最短行走時間最小。

\textbf{關鍵觀察。} 令 $\delta(v_a, v_b)$ 是把 $G$ 當成\textbf{無向圖}時 $v_a$ 到 $v_b$ 的最短距離。
\begin{itemize}
\item 不管怎麼定方向，定完之後的任何一條 $v_a$ 到 $v_b$ 的有向路徑，忽略方向後都是 $G$ 裡的一條路徑，
長度 $\ge \delta(v_a, v_b)$。 所以答案不可能比 $\delta(v_a, v_b)$ 小。
\item 取一條無向的最短路徑 $P$：$v_a = u_0, u_1, \ldots, u_k = v_b$。 最短路徑是簡單路徑（權重都是正的，
有重複頂點就能剪掉一段變得更短），每條邊只出現一次，把 $P$ 上的邊都定成 $u_{t-1} \to u_t$ 不會互相
衝突；其他的邊方向隨便定。 這樣 $v_a$ 沿 $P$ 走到 $v_b$ 花 $\delta(v_a, v_b)$。
\end{itemize}
所以最小值就是 $\delta(v_a, v_b)$，問題化成在無向圖上找一條最短路徑，再沿路定方向。
\begin{pseudo}{定方向讓 $v_a$ 到 $v_b$ 最快}
\Procedure{Orient}{$G, w, a, b$}
    \State 把 $G$ 當無向圖，從 $a$ 求最短路徑的前驅 $\pi$（(1) 用 BFS，(2) 用 Dijkstra）
    \State $v \gets b$
    \While{$v \ne a$}
        \State 把邊 $\{\pi[v], v\}$ 定成 $\pi[v] \to v$
        \State $v \gets \pi[v]$
    \EndWhile
    \State 其餘的邊方向任意定
\EndProcedure
\end{pseudo}

\textbf{(b)(1) 權重都相同。} 每條邊的時間都一樣，路徑時間和邊數成正比，最短路徑就是邊數最少的路徑，
用 BFS（資結書第~21.3~節，資結書定理~21.4）求出 $\pi$。 BFS 是 $O(n + m)$，沿 $\pi$ 走回去最多 $n - 1$ 步，
替其他邊定方向 $O(m)$，總共 $O(n + m)$，是線性時間。

\textbf{(b)(2) 權重不同。} 權重都是正的，用 Dijkstra 演算法（資結書第~25.2~節）在無向圖上求 $\pi$
（每條無向邊當成兩個方向的有向邊，正確性由資結書定理~25.6）。 用二元堆積時 $O((n + m) \lg n)$，用
Fibonacci 堆積時 $O(m + n \lg n)$；後面定方向的部分 $O(n + m)$。 總時間 $O((n + m) \lg n)$（或
$O(m + n \lg n)$）。 這裡不能用 BFS：邊數最少的路徑不一定時間最短，例如一段 $10$ 和兩段 $1$。
\end{pastsoln}
